\chapter{Counting}

\begin{center}\itshape
There are $n^{n-2}$ trees on $n$ labelled vertices.\\
\upshape(Borchardt, 1860; Cayley, 1889)
\end{center}

London, 1857. Arthur Cayley is a lawyer. He will spend fourteen years at the bar drafting conveyances, and in his spare
time he writes some two hundred papers of mathematics. That year he is composing differential operators, and he notices
that the terms they produce are organised like the branches of a tree: one trunk, which splits, whose branches split
again. He calls these diagrams \emph{trees}, and counts them: $1,1,2,4,9,20,48,\dots$ rooted trees with $1,2,3,\dots$ vertices.

Seventeen years later the trees come back from an unexpected direction. The chemists of the 1860s have learned that a
carbon atom makes four bonds and a hydrogen atom one. A molecule of formula $\mathrm C_n\mathrm H_{2n+2}$, a \emph{paraffin} (an
\emph{alkane}, today), is then a tree of $n$ carbon atoms, each with at most four carbon neighbours, the free bonds filled
with hydrogen. Different trees give different substances with the same formula: \emph{isomers}. Butane and isobutane are
both $\mathrm C_4\mathrm H_{10}$; there are three pentanes, and five hexanes (Figure~\ref{fig:hexanes}). Chemists were
synthesising them one by one and wanted to know how many to look for. In 1875 Cayley answered, up to thirteen carbon atoms,
with a method based on the \emph{centre} of a tree, an idea of Camille Jordan. Sixty years later, in Z\"urich, George P\'olya
turned his ad hoc computations into a single theorem, which also counts alcohols, benzene derivatives and graphs; we met it
in Tome~2, Chapter~3.

\begin{figure}[ht]
\centering
\fig{ch03-hexanes}
\caption{The five hexanes $\mathrm C_6\mathrm H_{14}$, drawn as their carbon skeletons (the hydrogens fill the free bonds). From
left to right: hexane, 2-methylpentane, 3-methylpentane, 2,2-dimethylbutane and 2,3-dimethylbutane. Section~\ref{sec:alkanes}
counts such skeletons for any number of carbons.}
\label{fig:hexanes}
\end{figure}

Counting is the oldest part of combinatorics and one of the most useful parts of computer science. It tells us how big a
search space is before we search it, how many inputs an algorithm must handle, how likely a hash collision is, how to
generate a random object uniformly, and how many genuinely different test cases a problem has once symmetric ones are
identified. This chapter is about the ideas that make counting \emph{exact}: double counting, the M\"obius function of a poset,
bijections (a tree is a word, Pr\"ufer's discovery), a little lemma on cyclic shifts that counts Catalan objects and ballots,
counting up to symmetry, and the $q$-analogues of binomial coefficients that count subspaces of vector spaces over finite
fields. Chapter~4 will add the machinery of generating functions.

\begin{think}
Draw all the trees on the four vertices $1,2,3,4$ (connected graphs with three edges). There are $16$ of them. If you erase
the labels, how many different shapes remain? Then guess how many trees there are on $5$ labelled vertices.
\end{think}

\section{Double counting}

The principle is simple: to count a set, count it twice, in two different ways, and equate. The set is usually a set of
pairs, counted once by grouping by the first coordinate and once by the second.

The first example is the \emph{handshake lemma}. In a graph with vertex set $V$ and edge set $E$, count the pairs $(v,e)$ with
$v$ an endpoint of $e$: grouping by $v$ gives $\sum_v\deg v$, grouping by $e$ gives $2|E|$. So $\sum_v\deg v=2|E|$, and in particular the
number of vertices of odd degree is even. At any party, the number of guests who shook an odd number of hands is even.

The second example is a loop that every programmer has written, the sieve:
\begin{center}
\texttt{for d in 1..n: for k in d, 2d, 3d, ..., up to n: visit(d, k)}
\end{center}
It visits each pair $(d,k)$ with $d\mid k\le n$ once. Grouping by $d$, there are $\lfloor n/d\rfloor$ multiples of $d$; grouping by $k$, there
are $\tau(k)$ divisors of $k$. Hence
\[
\sum_{k=1}^n\tau(k)=\sum_{d=1}^n\Big\lfloor\frac nd\Big\rfloor .
\]
Since $\frac nd-1<\lfloor\frac nd\rfloor\le\frac nd$ and $\ln(n+1)\le H_n=\sum_{d\le n}\frac1d\le1+\ln n$, the loop takes $n\ln n+O(n)$ steps, and an integer below $n$ has
$\ln n+O(1)$ divisors on average. Tome~1, Chapter~4 sharpens this with Dirichlet's hyperbola method, and Tome~1,
Proposition~4.5 proved Gauss' identity $\sum_{d\mid n}\varphi(d)=n$ by another double count. Burnside's lemma (Tome~2,
Theorem~3.6), the proof of the Erd\H os--Ko--Rado theorem in Chapter~2, and Pitman's proof of Cayley's formula
(Exercise~\ref{ex:pitman}) are double counts too.

\section{The M\"obius function of a poset}

Two formulas from earlier tomes look like cousins. The inclusion--exclusion principle (Tome~0, Theorem~8.2) recovers the
number of elements in none of the sets $A_1,\dots,A_n$ from the sizes of all intersections, with signs $(-1)^{|S|}$. M\"obius
inversion (Tome~1, Theorem~4.4) recovers $f$ from $g(n)=\sum_{d\mid n}f(d)$, with the signs $\mu(n/d)$. In 1964 Gian-Carlo Rota
showed that they are one theorem, about posets, and that every locally finite poset has its own M\"obius function. His paper,
\emph{On the foundations of combinatorial theory}, is often said to have turned combinatorics from a bag of tricks into a
subject.

A poset is \emph{locally finite} if every interval $[x,y]=\{z:x\le z\le y\}$ is finite. Its \emph{incidence algebra} is the set of functions
$f(x,y)$ with rational values defined for all pairs $x\le y$, with pointwise sum and the \emph{convolution}
\[
(f*g)(x,y)=\sum_{x\le z\le y}f(x,z)\,g(z,y),
\]
a finite sum. Two functions play a special role: $\delta(x,y)=1$ if $x=y$ and $0$ otherwise, and the \emph{zeta function} $\zeta(x,y)=1$
for all $x\le y$.

\begin{proposition}\label{prop:incidence}
The convolution is associative and $\delta$ is its identity. A function $f$ has a two-sided inverse if and only if $f(x,x)\ne0$
for all $x$; the inverse is then unique. If moreover $f$ has integer values and $f(x,x)=1$ for all $x$, its inverse has integer
values.
\end{proposition}
\begin{proof}
Both $((f*g)*h)(x,y)$ and $(f*(g*h))(x,y)$ equal the sum of $f(x,z)\,g(z,w)\,h(w,y)$ over the pairs $x\le z\le w\le y$, and $f*\delta=\delta*f=f$
is clear. If $f*g=\delta$ then $f(x,x)g(x,x)=1$, so $f(x,x)\neq0$. Conversely, define $g$ by induction on the size of the interval $[x,y]$:
\[
\begin{gathered}g(x,x)=\frac1{f(x,x)},\qq g(x,y)=-\frac1{f(y,y)}\sum_{x\le z<y}g(x,z)\,f(z,y);;(x<y),\end{gathered}
\]
which is exactly $(g*f)(x,y)=\delta(x,y)$. In the same way $h(x,x)=\frac1{f(x,x)}$ and $h(x,y)=-\frac1{f(x,x)}\sum_{x<z\le y}f(x,z)\,h(z,y)$ give
$f*h=\delta$. Then $g=g*(f*h)=(g*f)*h=h$, a two-sided inverse, and any other inverse $g'$ satisfies $g'=g'*(f*h)=h$. When
$f(x,x)=1$ and $f$ is integer valued, the recursion only adds and multiplies integers.
\end{proof}

\begin{definition}\label{def:mobius}
The \emph{M\"obius function} of a locally finite poset is $\mu=\zeta^{-1}$. Explicitly, $\mu(x,x)=1$ and
\[
\begin{aligned}{}&\mu(x,y)=-\sum_{x\le z<y}\mu(x,z)\eqbr=-\sum_{x<z\le y}\mu(z,y)\qquad\text{for }x<y.\end{aligned}
\]
\end{definition}

The two recursions are $\mu*\zeta=\delta$ and $\zeta*\mu=\delta$. Computing $\mu$ by hand is a walk up the Hasse diagram: write $1$ at $x$, and at
each element above it write minus the sum of what is written below it.

\begin{theorem}[M\"obius inversion, Rota]\label{thm:mobiusinv}
Let $P$ be a poset in which every set $\{x:x\le y\}$ is finite, and $f,g:P\to\Q$. Then
\[
\begin{gathered}g(y)=\sum_{x\le y}f(x)\ \text{ for all }y\qq\iff\qq f(y)=\sum_{x\le y}g(x)\,\mu(x,y)\ \text{ for all }y.\end{gathered}
\]
Dually, if every set $\{y:y\ge x\}$ is finite, $g(x)=\sum_{y\ge x}f(y)$ for all $x$ if and only if $f(x)=\sum_{y\ge x}\mu(x,y)\,g(y)$ for all $x$.
\end{theorem}
\begin{proof}
If $g(y)=\sum_{x\le y}f(x)$ for all $y$, then, exchanging two finite sums,
\[
\begin{aligned}{}&\sum_{x\le y}g(x)\mu(x,y)=\sum_{x\le y}\sum_{w\le x}f(w)\mu(x,y)\eqbr=\sum_{w\le y}f(w)\sum_{w\le x\le y}\zeta(w,x)\mu(x,y)\eqbr=\sum_{w\le y}f(w)\,\delta(w,y)=f(y).\end{aligned}
\]
Conversely, if $f(y)=\sum_{x\le y}g(x)\mu(x,y)$, then $\sum_{x\le y}f(x)=\sum_{w\le y}g(w)\sum_{w\le x\le y}\mu(w,x)\zeta(x,y)=\sum_{w\le y}g(w)\delta(w,y)=g(y)$.
The dual statement is the same computation with the products written in the other order.
\end{proof}

To use the theorem, we need $\mu$. Three posets cover most applications, and a product rule computes all three.

\begin{proposition}\label{prop:mobiusproduct}
If $P$ and $Q$ are locally finite, the M\"obius function of the product order on $P\times Q$ is
$\mu((x,x'),(y,y'))=\mu_P(x,y)\,\mu_Q(x',y')$. Consequently:
\begin{enumerate}
\item on the chain $0<1<2<\cdots$, $\mu(i,i)=1$, $\mu(i,i+1)=-1$ and $\mu(i,j)=0$ for $j\ge i+2$;
\item on the Boolean lattice of subsets of a finite set, $\mu(S,T)=(-1)^{|T\setminus S|}$ for $S\subseteq T$;
\item on the positive integers ordered by divisibility, $\mu(d,e)=\mu(e/d)$ for $d\mid e$, the classical M\"obius function of Tome~1.
\end{enumerate}
\end{proposition}
\begin{proof}
Let $\nu((x,x'),(y,y'))=\mu_P(x,y)\mu_Q(x',y')$. The interval from $(x,x')$ to $(y,y')$ is $[x,y]\times[x',y']$, so
\[
\begin{aligned}{}&\sum_{(x,x')\le(z,z')\le(y,y')}\nu((x,x'),(z,z'))\eqbr=\Big(\sum_{x\le z\le y}\mu_P(x,z)\Big)\Big(\sum_{x'\le z'\le y'}\mu_Q(x',z')\Big)\eqbr=\delta(x,y)\,\delta(x',y'),\end{aligned}
\]
that is, $\nu*\zeta=\delta$, and $\nu=\mu$ by uniqueness of inverses. (1) The recursion gives $\mu(i,i+1)=-1$, then $\mu(i,i+2)=-(1-1)=0$, and
by induction $\mu(i,j)=-(1-1+0+\dots+0)=0$ for $j\ge i+2$. (2) The subsets of $\{1,\dots,n\}$ are the product of $n$ chains $\{0<1\}$ (a
subset is its indicator vector), so $\mu(S,T)$ is a product of $|S|$ factors $\mu(1,1)=1$, $|T\setminus S|$ factors $\mu(0,1)=-1$ and factors
$\mu(0,0)=1$. (3) The divisors of $e$ form the product of the chains $\{0,1,\dots,a_p\}$ of exponents, $e=\prod p^{a_p}$. So $\mu(d,e)$ is
$0$ if some prime divides $e/d$ twice, and $(-1)^r$ if $e/d$ is a product of $r$ distinct primes: this is $\mu(e/d)$.
\end{proof}

Case~(3) turns Theorem~\ref{thm:mobiusinv} into the M\"obius inversion of Tome~1. Case~(2) gives inclusion--exclusion. Let
$A_1,\dots,A_n$ be subsets of a finite set $\Omega$, and for $S\subseteq\{1,\dots,n\}$ let $f(S)$ be the number of elements of $\Omega$ that lie in $A_i$
exactly for $i\in S$, and $g(S)=|\bigcap_{i\in S}A_i|$ (with $g(\emptyset)=|\Omega|$). An element counted by $g(S)$ lies in the $A_i$ for $i$ in some set
$T\supseteq S$, so $g(S)=\sum_{T\supseteq S}f(T)$. The dual inversion gives $f(S)=\sum_{T\supseteq S}(-1)^{|T\setminus S|}g(T)$, and at $S=\emptyset$:
\[
\Big|\Omega\setminus\bigcup_iA_i\Big|=\sum_{T\subseteq\{1,\dots,n\}}(-1)^{|T|}\Big|\bigcap_{i\in T}A_i\Big| .
\]
One proof, two classical formulas, and infinitely many new ones: every poset now has an inversion formula.

\begin{think}
Compute $\mu(\hat0,\hat1)$ for the poset with a bottom $\hat0$, a top $\hat1$ and $k$ incomparable elements in between. For which $k$ is it zero?
\end{think}

A last formula shows what $\mu$ measures. Call a \emph{chain of length $k$} from $x$ to $y$ a sequence $x=x_0<x_1<\dots<x_k=y$.

\begin{proposition}[Philip Hall, 1936]\label{prop:hall}
In a locally finite poset, $\mu(x,y)=\sum_{k\ge0}(-1)^kc_k(x,y)$, where $c_k(x,y)$ is the number of chains of length $k$ from $x$ to $y$.
\end{proposition}
\begin{proof}
Let $\eta=\zeta-\delta$, so $\eta(x,y)=1$ if $x<y$ and $0$ otherwise. Expanding the convolution, $\eta^{*k}(x,y)$ is the number of sequences
$x=x_0<x_1<\dots<x_k=y$, that is, $c_k(x,y)$; it vanishes for $k\ge|[x,y]|$, since a chain of length $k$ has $k+1$ distinct elements. So
$\nu(x,y)=\sum_{k\ge0}(-1)^k\eta^{*k}(x,y)$ is a finite sum for each pair. The value of $\zeta*\nu$ at $(x,y)$ only involves the interval
$[x,y]$, on which $\nu$ agrees with $\sum_{k=0}^N(-\eta)^{*k}$ for $N=|[x,y]|$, and
\[
\begin{aligned}{}&\zeta*\sum_{k=0}^N(-\eta)^{*k}=(\delta+\eta)*\sum_{k=0}^N(-\eta)^{*k}\eqbr=\delta-(-\eta)^{*(N+1)},\end{aligned}
\]
a telescoping sum, whose last term vanishes on $[x,y]$. So $\zeta*\nu=\delta$ and $\nu=\mu$.
\end{proof}

\begin{remark}
The chains of the open interval $(x,y)$ form a simplicial complex, and Hall's formula says that $\mu(x,y)$ is its reduced Euler
characteristic. This is the start of \emph{topological combinatorics}: the M\"obius function of the lattice of partitions of a set,
$\mu(\hat0,\hat1)=(-1)^{n-1}(n-1)!$ (Exercise~\ref{ex:partitionlattice}), comes from a complex that is, up to homotopy, a bouquet of $(n-1)!$
spheres. See Stanley's \emph{Enumerative Combinatorics}, Volume~1, Chapter~3.
\end{remark}

\subsection*{Primitive words}

A word of length $n$ over an alphabet of $q$ letters is \emph{primitive} if it is not a power $u^m=uu\cdots u$ with $m\ge2$: \texttt{abab} is not,
\texttt{abba} is. The \emph{rotation} by $r$ of $w=w_0w_1\cdots w_{n-1}$ is $w_rw_{r+1}\cdots w_{r-1}$ (indices mod $n$).

\begin{lemma}\label{lem:primitive}
Every non-empty word $w$ is $u^{n/d}$ for exactly one primitive word $u$, whose length $d$ divides $n=|w|$.
\end{lemma}
\begin{proof}
For $d\mid n$, $w=u^{n/d}$ with $|u|=d$ if and only if $w_i=w_{i+d}$ for all $i$ (mod $n$), that is, the rotation by $d$ fixes $w$. The rotations that
fix $w$ form a subgroup of $\Z/n\Z$, hence (Tome~2, Chapter~1) the multiples of some $d\mid n$, its least positive element. So $w=u^{n/d}$
where $u$ is the prefix of length $d$, and $u$ is primitive: if $u=v^r$ with $r\ge2$, the rotation by $|v|<d$ would fix $w$. If also $w=v^{n/e}$
with $v$ primitive, the rotation by $e$ fixes $w$, so $d\mid e$ and $v=u^{e/d}$; as $v$ is primitive, $e=d$ and $v=u$.
\end{proof}

So if $M(d)$ is the number of primitive words of length $d$, then $q^n=\sum_{d\mid n}M(d)$, and Theorem~\ref{thm:mobiusinv} on the
divisor poset gives
\[
M(n)=\sum_{d\mid n}\mu\Big(\frac nd\Big)\,q^d .
\]
For $q=2$: $M(1)=2$, $M(2)=2$, $M(3)=6$, $M(4)=12$, $M(6)=54$. The rotations of a primitive word are $n$ distinct primitive words, so
primitive words come in classes of $n$, and there are $M(n)/n$ classes. This is \emph{exactly} the number of monic irreducible
polynomials of degree $n$ over $\mathbb F_q$ (Tome~2, Theorem~9.10), and also the dimension of the degree-$n$ part of the free Lie
algebra on $q$ generators (Witt's formula, 1937). Exercise~\ref{ex:lyndon} picks one word in each class, the Lyndon word, which
is how these coincidences are turned into bijections.

\section{Trees are words}

We need a few facts about trees here; Chapter~5 studies graphs in general. A \emph{graph} is a set $V$ of vertices with a set $E$ of
\emph{edges}, which are pairs of distinct vertices. A \emph{path} is a sequence of distinct vertices with consecutive ones adjacent, a \emph{cycle}
is a path $v_0\cdots v_k$ with $k\ge2$ and $v_kv_0$ an edge, and the graph is \emph{connected} if any two vertices are joined by a path. A \emph{tree} is a
connected graph without cycles, and a \emph{leaf} is a vertex of degree $1$.

\begin{lemma}\label{lem:trees}
Let $T$ be a tree with $n\ge2$ vertices.
\begin{enumerate}
\item $T$ has at least two leaves.
\item If $v$ is a leaf, $T-v$ is a tree.
\item $T$ has $n-1$ edges.
\item Adding a new vertex joined to one vertex of $T$ gives a tree.
\end{enumerate}
\end{lemma}
\begin{proof}
(1) As $T$ is connected with $n\ge2$ vertices, it has an edge. Take a path $v_0v_1\cdots v_k$ with $k\ge1$ as long as possible. A neighbour $u\ne v_1$
of $v_0$ is either some $v_i$ with $i\ge2$, and $v_0v_1\cdots v_i$ is a cycle, or not on the path, and $uv_0\cdots v_k$ is longer. Both are impossible, so
$v_0$ is a leaf, and so is $v_k$. (2) $T-v$ has no cycle. Two vertices $x,y\ne v$ are joined in $T$ by a path, whose interior vertices have
two neighbours on it, so $v$ is not among them, and the path lives in $T-v$. (3) follows from (1) and (2) by induction on $n$. (4) The
new vertex has degree $1$, and a cycle has two edges at each of its vertices, so a cycle of the new graph would be a cycle of $T$.
\end{proof}

Now take a tree $T$ whose $n\ge2$ vertices are labelled by distinct integers, so that ``the smallest leaf'' makes sense. Its
\emph{Pr\"ufer code} $P(T)$ is the word of length $n-2$ defined by: $P(T)$ is empty if $n=2$; otherwise, if $\ell$ is the smallest leaf and $a$ its
neighbour, $P(T)=a\,P(T-\ell)$. In words: pluck the smallest leaf, write down its neighbour, repeat until two vertices remain
(Figure~\ref{fig:prufer}).

\begin{figure}[ht]
\centering
\fig{ch03-prufer}
\caption{Computing a Pr\"ufer code. The leaves are plucked in the order $1,2,3,4,5$ (the circled numbers give the step), and
their neighbours are written down: the code is $(4,4,5,5,6)$. Each vertex appears its degree minus one times: $4$ and $5$ twice,
$6$ once, the leaves $1,2,3,7$ never.}
\label{fig:prufer}
\end{figure}

\begin{lemma}\label{lem:pruferdeg}
Every vertex $v$ of $T$ appears exactly $\deg_T(v)-1$ times in $P(T)$. In particular, the leaves of $T$ are the vertices absent from $P(T)$.
\end{lemma}
\begin{proof}
By induction on $n$. For $n=2$ both vertices have degree $1$ and the code is empty. For $n\ge3$, let $T'=T-\ell$. In $T'$, the vertex $a$ has
lost one edge and the other degrees are unchanged. By induction $v\ne\ell$ appears $\deg_{T'}(v)-1$ times in $P(T')$; adding the first
letter $a$ gives $\deg_T(a)-1$ occurrences of $a$ and $\deg_T(v)-1$ of the other $v\ne\ell$. Finally $\ell$ is not a vertex of $T'$, so it occurs
$0=\deg_T(\ell)-1$ times.
\end{proof}

\begin{theorem}[Pr\"ufer, 1918; Cayley's formula]\label{thm:prufer}
Let $V$ be a set of $n\ge2$ integers. The Pr\"ufer code is a bijection from the trees with vertex set $V$ to the words $V^{n-2}$.
In particular there are $n^{n-2}$ trees on $n$ labelled vertices.
\end{theorem}
\begin{proof}
By induction on $n$; for $n=2$ there is one tree and one empty word. Let $n\ge3$ and $a=(a_1,\dots,a_{n-2})\in V^{n-2}$. Let $\ell$ be the smallest
element of $V$ that does not occur in $a$ (there is one, as $a$ has $n-2<n$ letters).

\emph{At most one tree has code $a$.} If $P(T)=a$, then by Lemma~\ref{lem:pruferdeg} the leaves of $T$ are the elements of $V$ absent from $a$, so the
smallest leaf of $T$ is $\ell$, its neighbour is $a_1$, and $P(T-\ell)=(a_2,\dots,a_{n-2})$. By Lemma~\ref{lem:trees}, $T-\ell$ is a tree on
$V\setminus\{\ell\}$, which by induction is determined by its code; and $T$ is $T-\ell$ plus the edge $\ell a_1$.

\emph{Some tree has code $a$.} The letters $a_2,\dots,a_{n-2}$ lie in $V\setminus\{\ell\}$, so by induction there is a tree $T'$ on $V\setminus\{\ell\}$ with code
$(a_2,\dots,a_{n-2})$. Let $T$ be $T'$ plus the edge $\ell a_1$ (note $a_1\ne\ell$), a tree by Lemma~\ref{lem:trees}. For $v\ne\ell$, Lemma~\ref{lem:pruferdeg}
applied to $T'$ gives $\deg_T(v)=1+(\text{occurrences of }v\text{ in }a_2\cdots a_{n-2})+[v=a_1]$, which is $1$ exactly when $v$ does not occur in $a$. So
the leaves of $T$ are $\ell$ and the elements of $V\setminus\{\ell\}$ absent from $a$, and $\ell$ is the smallest of them. Hence $P(T)=a_1P(T')=a$.
\end{proof}

Lemma~\ref{lem:pruferdeg} says more: the code of a tree knows its degrees. Counting words with prescribed numbers of
occurrences of each letter is the multinomial theorem.

\begin{corollary}\label{cor:degrees}
Let $d_1,\dots,d_n\ge1$ with $d_1+\dots+d_n=2n-2$. The number of trees on $\{1,\dots,n\}$ in which vertex $i$ has degree $d_i$ for all $i$ is
\[
\begin{aligned}{}&\binom{n-2}{d_1-1,\ d_2-1,\ \dots,\ d_n-1}\eqbr=\frac{(n-2)!}{(d_1-1)!\cdots(d_n-1)!}.\end{aligned}
\]
Equivalently, $\displaystyle\sum_T x_1^{\deg_T(1)}\cdots x_n^{\deg_T(n)}=x_1x_2\cdots x_n\,(x_1+\dots+x_n)^{n-2}$, the sum over all trees on $\{1,\dots,n\}$.
\end{corollary}
\begin{proof}
By Theorem~\ref{thm:prufer} and Lemma~\ref{lem:pruferdeg}, these trees correspond to the words of length $n-2$ in which $i$ occurs $d_i-1$
times, and choosing the positions of each letter gives the multinomial coefficient. For the identity, the word $a$ of a tree
$T$ satisfies $\prod_ix_i^{\deg_T(i)}=x_1\cdots x_n\,x_{a_1}\cdots x_{a_{n-2}}$, and the sum of $x_{a_1}\cdots x_{a_{n-2}}$ over all words is $(x_1+\dots+x_n)^{n-2}$.
\end{proof}

For $n=4$: degrees $(3,1,1,1)$ give $\frac{2!}{2!}=1$ tree, the star centred at $1$; degrees $(2,2,1,1)$ give $\frac{2!}{1!\,1!}=2$ trees, the paths
$3\,1\,2\,4$ and $4\,1\,2\,3$. Over all choices of centres, $4$ stars and $6\times2=12$ paths: the $16=4^2$ trees of the opening, with two
shapes.

\begin{beauty}
A tree, a two-dimensional branching object, is the same thing as a word of $n-2$ letters, and the degree of each vertex can be
read from the word by counting. This one bijection proves Cayley's formula, counts trees by degrees, and gives an algorithm: to
draw a uniformly random tree on $n$ vertices, draw $n-2$ independent uniform letters and decode, in time $O(n\log n)$ with a heap.
The companion notebook does exactly this, and checks with it a fact of Exercise~\ref{ex:leaf}: about $n/e$ of the vertices of a random
tree are leaves.
\end{beauty}

\begin{remark}
Cayley's formula has an unusual number of beautiful proofs. Borchardt's original one (1860) used determinants, and is the
case of complete graphs of the matrix-tree theorem of Chapter~9. Joyal (1981) matched doubly rooted trees with maps from
$\{1,\dots,n\}$ to itself; Pitman (1999) found the double count of Exercise~\ref{ex:pitman}. Aigner and Ziegler's \emph{Proofs from THE BOOK}
gives four of them. Tome~3, Chapter~12 used the formula to decide when a random graph is connected.
\end{remark}

\section{The cycle lemma}

In 1887 Joseph Bertrand asked: in an election where candidate A receives $a$ votes and B receives $b<a$, counted in random order,
what is the probability that A is ahead throughout the count? The answer, $\frac{a-b}{a+b}$, was proved in Tome~3, Theorem~1.12 by
D\'esir\'e Andr\'e's reflection trick. In 1947 Aryeh Dvoretzky and Theodore Motzkin found a proof that needs no reflection at all,
only rotation, and it counts much more than ballots.

\begin{theorem}[cycle lemma, Dvoretzky and Motzkin]\label{thm:cyclelemma}
Let $x_1,\dots,x_m$ be integers, each at most $1$, with sum $k\ge1$. Among the $m$ cyclic shifts $x_{i+1},x_{i+2},\dots,x_{i+m}$ ($0\le i<m$, indices mod $m$),
exactly $k$ have all their partial sums positive.
\end{theorem}
\begin{proof}
Extend the sequence periodically, $x_{j+m}=x_j$ for all $j\in\Z$, and let $S_j$ be its walk: $S_0=0$ and $S_j-S_{j-1}=x_j$ for all $j\in\Z$. Then
$S_{j+m}=S_j+k$, so $S_j\to+\infty$ as $j\to+\infty$ and $S_j\to-\infty$ as $j\to-\infty$; and the walk goes up by at most $1$ at each step. The shift starting
after $i$ has all partial sums positive if and only if $S_j>S_i$ for $i<j\le i+m$; call $i$ \emph{good} if $S_j>S_i$ for \emph{all} $j>i$. The two conditions
are the same, because $S_j=S_{j-m}+k>S_{j-m}$ propagates the inequality from one period to the next.

Since the walk climbs by steps of at most $1$ from $-\infty$ to $+\infty$, it visits every level $v\in\Z$, and finitely often; let $t(v)$ be the last
time it is at level $v$. After $t(v)$ the walk never goes below $v$ either, since it would have to climb back through $v$. So $t(v)$ is
good, and conversely a good $i$ is the last visit $t(S_i)$ to its level. Thus $v\mapsto t(v)$ is a bijection from $\Z$ to the good times, and
it is increasing: if $v<w$, the walk is above $w>v$ after $t(w)$, so $t(v)<t(w)$. Finally $t(v+k)=t(v)+m$, by periodicity. The good
times in the window $t(0)\le i<t(k)=t(0)+m$ are therefore $t(0),\dots,t(k-1)$: exactly $k$ of them. As the set of good times is invariant
under $i\mapsto i+m$, every window of $m$ consecutive integers, in particular $\{0,\dots,m-1\}$, contains exactly $k$.
\end{proof}

\begin{figure}[ht]
\centering
\fig{ch03-cyclelemma}
\caption{The cycle lemma for the sequence $(1,-2,1,1,-1,1,1)$, of sum $k=2$, extended periodically. A time is good when the
walk stays strictly above its current level forever after (dashed lines). The good times are the last visits to each
level, and there are exactly $2$ per period: the shifts starting after positions $2$ and $5$, that is $(1,1,-1,1,1,1,-2)$ and
$(1,1,1,-2,1,1,-1)$, have all partial sums positive.}
\label{fig:cyclelemma}
\end{figure}

\begin{corollary}[ballots and Catalan numbers]\label{cor:ballot}
Let $a>b\ge0$. The number of sequences of $a$ terms $+1$ and $b$ terms $-1$ whose partial sums are all positive is
$\frac{a-b}{a+b}\binom{a+b}a$. In particular, the number of \emph{Dyck words} of length $2n$, sequences of $n$ terms $+1$ and $n$ terms $-1$
whose partial sums are all $\ge0$, is the Catalan number
\[
C_n=\frac1{2n+1}\binom{2n+1}n=\frac1{n+1}\binom{2n}n .
\]
\end{corollary}
\begin{proof}
Let $m=a+b$ and $k=a-b$. Count the pairs $(s,i)$ where $s$ is one of the $\binom ma$ sequences and $0\le i<m$ is a position such that the
shift of $s$ starting after $i$ has positive partial sums. By the cycle lemma each $s$ has $k$ such positions: $k\binom ma$ pairs. On the
other hand, each sequence $g$ with positive partial sums arises from exactly $m$ pairs, namely $s$ the rotation of $g$ by $-i$, with this
$i$, for $i=0,\dots,m-1$. Hence the number of such $g$ is $\frac km\binom ma$. For Dyck words, $w$ has partial sums $\ge0$ if and only if $(+1)w$ has
partial sums $\ge1$; these are the sequences of $n+1$ terms $+1$ and $n$ terms $-1$ with positive partial sums, and there are
$\frac1{2n+1}\binom{2n+1}{n+1}=\frac1{n+1}\binom{2n}n$ of them.
\end{proof}

The proof divides by $m$, and the division is exact for a reason that a computer scientist loves: the $m$ rotations of the
\emph{pairs} $(s,i)$ always form a full orbit, even when $s$ has symmetries, because the position $i$ breaks them.

The Catalan numbers $1,1,2,5,14,42,132,\dots$ count more than two hundred families of objects (Stanley's \emph{Catalan Numbers}
lists 214). Here is the one that matters most for computing. A \emph{plane tree} is a rooted tree in which the children of each vertex
are ordered, left to right: the shape of a parse tree, an XML document, or a recursive call. Visit its vertices in \emph{preorder}
(the root, then the subtrees from left to right, each in preorder), and record the number of children $c_1,\dots,c_m$ of each.

\begin{proposition}[\L ukasiewicz words]\label{prop:lukasiewicz}
A sequence $(c_1,\dots,c_m)$ of non-negative integers is the preorder sequence of numbers of children of a plane tree if and
only if $h_j=1+\sum_{i\le j}(c_i-1)$ satisfies $h_j\ge1$ for $j<m$ and $h_m=0$; and the tree is then unique. Consequently the number of plane
trees with $m$ vertices whose numbers of children form a given multiset with sum $m-1$ is $\frac1m$ times the number of
arrangements of the multiset; in particular there are $C_{m-1}$ plane trees with $m$ vertices.
\end{proposition}
\begin{proof}
More generally, call a sequence a \emph{code of $r$ trees} if $h_j=r+\sum_{i\le j}(c_i-1)$ satisfies $h_j\ge1$ for $j<m$ and $h_m=0$; think of $h_j$ as the
number of subtrees still to be read after $j$ vertices. We show by induction on $m$ that the codes of $r$ trees are exactly the
concatenations of the preorder sequences of $r$ plane trees, each obtained in exactly one way. The preorder sequence of a tree
is $c_1$ followed by the sequences of its $c_1$ subtrees, so the claim for $r$ trees with $m$ vertices in all reduces to the claim for $c_1$
trees, then $r-1$ trees, with fewer vertices. For the converse, $h$ starts at $r$ and goes down by at most $1$ at each step, so a code of $r$
trees first reaches $r-1$ at some time $j_1$; the part up to $j_1$ is a code of one tree (shifted by $r-1$), the rest a code of $r-1$ trees; and
a code of one tree, read after its first term $c_1$, is a code of $c_1$ trees. The induction hypothesis applies to each piece, and the
cutting points are forced, which gives uniqueness.

For the count, set $x_i=1-c_i\le1$, of sum $m-(m-1)=1$. The condition on $h$ says that every \emph{suffix} $x_{j+1}+\dots+x_m=1-\sum_{i\le j}x_i=h_j$ for
$j<m$ is positive, and the full sum is $1$: the reversed sequence $(x_m,\dots,x_1)$ has positive partial sums. By the cycle lemma with $k=1$
and the pair count of Corollary~\ref{cor:ballot}, a fraction exactly $\frac1m$ of the arrangements qualifies. For all plane trees with $m$
vertices, the sequences of $m$ non-negative integers with sum $m-1$ number $\binom{2m-2}{m-1}$ (stars and bars), and $\frac1m\binom{2m-2}{m-1}=C_{m-1}$.
\end{proof}

This is how a stack-based parser checks a prefix expression such as \texttt{+ * a b c}: arities $2,2,0,0,0$, counter
$1,2,3,2,1,0$, accepted. Exercise~\ref{ex:fusscatalan} counts $r$-ary trees in the same way. In 1960 George Raney rediscovered the
cycle lemma and used it to prove the Lagrange inversion formula, which Chapter~4 will prove by another route.

\section{Counting up to symmetry}\label{sec:alkanes}

\begin{reminder}[Tome~2, Theorems~3.6 and~3.7]
Let a finite group $G$ act on a finite set $X$. \emph{Burnside}: the number of orbits is $\frac1{|G|}\sum_g|X^g|$, the average number of
fixed points. \emph{P\'olya}: if $g$ has $c_j(g)$ cycles of length $j$ on $X$, and the colours $c$ have weights $w_c$, the colourings of $X$ up to the
action of $G$, each weighted by the product of the weights of its colours, have total weight $Z_G(t_1,t_2,\dots)$ evaluated at
$t_j=\sum_cw_c^{\,j}$, where $Z_G=\frac1{|G|}\sum_g\prod_jt_j^{c_j(g)}$ is the \emph{cycle index}.
\end{reminder}

Tome~2 counted necklaces. Here we count graphs and trees, and finally answer Cayley's question about alkanes.

\subsection*{Graphs up to isomorphism}

A graph on the vertex set $\{1,\dots,n\}$ is a colouring of the $\binom n2$ pairs in two colours, ``edge'' and ``no edge''. Two graphs are
isomorphic if a permutation of the vertices maps one to the other, so the unlabelled graphs are the orbits of the symmetric
group $S_n$ acting on colourings of pairs. To apply P\'olya we need the cycles of a permutation $\sigma$ on \emph{pairs}. For $n=4$ (six
pairs), by conjugacy class of $S_4$:
\begin{itemize}
\item the identity fixes the $6$ pairs: $t_1^6$;
\item a transposition, $(12)$, fixes $\{1,2\}$ and $\{3,4\}$ and swaps $\{1,3\}\leftrightarrow\{2,3\}$ and $\{1,4\}\leftrightarrow\{2,4\}$: $t_1^2t_2^2$, $6$ of them;
\item a double transposition, $(12)(34)$, fixes $\{1,2\}$, $\{3,4\}$ and swaps $\{1,3\}\leftrightarrow\{2,4\}$, $\{1,4\}\leftrightarrow\{2,3\}$: $t_1^2t_2^2$, $3$ of them;
\item a $3$-cycle, $(123)$, rotates $\{1,2\}\to\{2,3\}\to\{1,3\}$ and $\{1,4\}\to\{2,4\}\to\{3,4\}$: $t_3^2$, $8$ of them;
\item a $4$-cycle, $(1234)$, rotates the four sides of the square and swaps the two diagonals: $t_2t_4$, $6$ of them.
\end{itemize}
So $Z=\frac1{24}\big(t_1^6+9t_1^2t_2^2+8t_3^2+6t_2t_4\big)$. With $t_j=1+x^j$ (an edge has weight $x$),
\[
\begin{aligned}{}&\tfrac1{24}\big((1+x)^6+9(1+x)^2(1+x^2)^2\eqbr{}+8(1+x^3)^2+6(1+x^2)(1+x^4)\big)\eqbr=1+x+2x^2+3x^3+2x^4+x^5+x^6:\end{aligned}
\]
there are $11$ graphs on $4$ vertices, of which $3$ have three edges (Figure~\ref{fig:graphs4}). With $5$ vertices there are $34$
(Exercise~\ref{ex:graphs5}), then $156$, $1044$, $12346$, \dots{} Since a single graph has at most $n!$ labellings, there are at least
$2^{\binom n2}/n!$ unlabelled graphs; P\'olya, and later Erd\H os and R\'enyi, showed that this is asymptotically exact, because almost all
graphs have no symmetry at all.

\begin{figure}[ht]
\centering
\fig{ch03-graphs4}
\caption{The $11$ graphs on $4$ unlabelled vertices, sorted by number of edges: $1,1,2,3,2,1,1$, the coefficients of P\'olya's
polynomial. Complementing (swapping edges and non-edges) reverses the list.}
\label{fig:graphs4}
\end{figure}

\subsection*{Multisets and rooted trees}

Most counting problems up to symmetry have the same shape: an object is made of several sub-objects whose order does not
matter. Let $\mathcal A$ be a class of objects, each with a size $\ge1$, with $a_n$ objects of size $n$ (finitely many), and $A(x)=\sum_na_nx^n$. A
\emph{multiset of $k$ elements} of $\mathcal A$ is a $k$-tuple up to permutation, that is, an orbit of $S_k$ acting on $\mathcal A^k$.

\begin{proposition}\label{prop:multisets}
The generating function, by total size, of the multisets of $k$ elements of $\mathcal A$ is $Z_{S_k}(A(x),A(x^2),A(x^3),\dots)$, where $Z_{S_k}$ is the
cycle index of $S_k$ acting on $\{1,\dots,k\}$. The generating function of all finite multisets is $\prod_{n\ge1}(1-x^n)^{-a_n}$.
\end{proposition}
\begin{proof}
A $k$-tuple is a colouring of $k$ slots, the colours being the objects of $\mathcal A$ and the weight of a colour $x^{\text{size}}$. Then $\sum_cw_c^{\,j}=A(x^j)$,
and P\'olya's theorem gives the first formula. There are infinitely many colours, but the coefficient of $x^N$ only involves
objects of size $\le N$, finitely many, so we may apply the theorem to them. For the second, a multiset is determined by the
multiplicity $m_o\ge0$ of each object $o$, and contributes $\prod_ox^{m_o|o|}$; summing over the choices independently gives
$\prod_o(1+x^{|o|}+x^{2|o|}+\cdots)=\prod_n(1-x^n)^{-a_n}$, a product in which only finitely many factors affect each coefficient.
\end{proof}

The small cycle indices, from the conjugacy classes of $S_k$:
\[
\begin{gathered}Z_{S_2}=\frac{t_1^2+t_2}2,\qquad Z_{S_3}=\frac{t_1^3+3t_1t_2+2t_3}6,\qq Z_{S_4}=\frac{t_1^4+6t_1^2t_2+3t_2^2+8t_1t_3+6t_4}{24}.\end{gathered}
\]

A rooted tree is a root together with the multiset of subtrees hanging from its children. Two rooted trees are isomorphic if and
only if these multisets agree up to isomorphism: an isomorphism sends the root to the root and children to children, so it
matches the subtrees; conversely, isomorphisms between matched subtrees glue together. So if $t_n$ is the number of rooted trees
with $n$ vertices up to isomorphism and $T(x)=\sum t_nx^n$, Proposition~\ref{prop:multisets} gives Cayley's equation of 1857:
\[
\begin{aligned}{}&T(x)=x\prod_{n\ge1}(1-x^n)^{-t_n}\eqbr=x+x^2+2x^3+4x^4+9x^5+20x^6+\cdots\end{aligned}
\]
It determines the $t_n$ one after the other, since the coefficient of $x^{n+1}$ on the right only involves $t_1,\dots,t_n$; Exercise~\ref{ex:rootedrec}
turns it into a fast recursion.

\subsection*{Alcohols and alkanes}

Now the chemistry. An \emph{alcohol} $\mathrm C_n\mathrm H_{2n+1}\mathrm{OH}$ is an alkane skeleton with one hydrogen replaced by a hydroxyl group: a tree
of $n$ carbons of degrees $\le4$ with a marked carbon that has at most three carbon neighbours. Rooted at the marked carbon, every
vertex has at most three children. Let $r_n$ be the number of such \emph{ternary rooted trees} with $n$ vertices, up to isomorphism, with
$r_0=1$ for the empty tree (which a chemist calls a hydrogen atom), and $R(x)=\sum_{n\ge0}r_nx^n$. A non-empty ternary tree is a root with
three slots, each holding a possibly empty ternary tree, up to permutation of the slots, so
\[
\begin{aligned}{}&R(x)=1+x\,Z_{S_3}\big(R(x),R(x^2),R(x^3)\big)\eqbr=1+x\,\frac{R(x)^3+3R(x)R(x^2)+2R(x^3)}6 .\end{aligned}
\]
This gives $1,1,1,2,4,8,17,39,89,211,507,\dots$: one methanol, one ethanol, two propanols, four butanols, eight pentanols, in
agreement with the chemists (Henze and Blair, 1931).

For the alkanes themselves the tree has no root, and we need a canonical vertex to root it at. For a vertex $v$ of a tree $T$, the
\emph{branches} at $v$ are the components of $T-v$.

\begin{lemma}[Jordan, 1869]\label{lem:centroid}
Let $T$ be a tree with $n$ vertices. Exactly one of the following holds: (i) exactly one vertex has all its branches of size
$<\frac n2$; (ii) $n$ is even and exactly one edge splits $T$ into two components of size $\frac n2$.
\end{lemma}
\begin{proof}
Call $v$ \emph{central} if all its branches have size $\le\frac n2$. Some vertex is central: start anywhere, and while the current vertex $v$ has a
branch $B$ of size $>\frac n2$, move to the neighbour $u$ of $v$ in $B$. The branches at $u$ are the component of $v$, of size $n-|B|<\frac n2$, and
the branches of $u$ inside $B$, of size $<|B|$. So the size of the large branch decreases strictly, and the walk stops at a central vertex.

Let $u\ne w$ be central, $A$ the branch at $u$ containing $w$, and $B$ the branch at $w$ containing $u$. Every vertex is in $A$ or in $B$:
if $x\notin A$, the path from $x$ to $w$ goes through $u$, and its part from $x$ to $u$ avoids $w$, so $x\in B$. Since $|A|,|B|\le\frac n2$, we get
$|A|=|B|=\frac n2$ and $A\cap B=\emptyset$. If $u$ and $w$ were not adjacent, an inner vertex of the path from $u$ to $w$ would be in $A\cap B$. So $u$ and $w$
are adjacent, and the edge $uw$ splits $T$ into $B$ and $A$, of size $\frac n2$ each: a \emph{halving edge}. Conversely both ends of a halving edge
are central. Two distinct halving edges would have at least three distinct ends, central and pairwise adjacent, a triangle,
impossible in a tree. Finally, a central vertex $v$ with a branch $C$ of size exactly $\frac n2$ is the end of a halving edge, the edge from
$v$ into $C$.

So either there is a halving edge, unique, whose two ends are the only central vertices (case (ii)), or there is none, and
then there is exactly one central vertex, and all its branches have size $<\frac n2$ (case (i)).
\end{proof}

\begin{theorem}[counting alkanes]\label{thm:alkanes}
Let $\rho(x)=\sum_{m<n/2}r_mx^m$, and $\varepsilon_n=\binom{r_{n/2}+1}2$ if $n$ is even, $\varepsilon_n=0$ if $n$ is odd. The number of alkanes $\mathrm C_n\mathrm H_{2n+2}$,
that is, of trees with $n$ vertices of degrees $\le4$ up to isomorphism, is
\[
\begin{aligned}{}&a_n=[x^{n-1}]\,Z_{S_4}\big(\rho(x),\rho(x^2),\rho(x^3),\rho(x^4)\big)\eqbr{}+\varepsilon_n.\end{aligned}
\]
\end{theorem}
\begin{proof}
An isomorphism of trees maps branches to branches of the same size, hence the vertex of case (i) of Lemma~\ref{lem:centroid} to the
vertex of case (i), and the halving edge to the halving edge. Count the two cases separately.

(i) Root the tree at its centroid $c$. Each branch at $c$, rooted at the neighbour of $c$, is a ternary rooted tree (its vertices have
at most $4$ neighbours, one of which is the parent), with fewer than $\frac n2$ vertices; there are at most $4$ branches, of total size $n-1$.
Conversely, a root with at most $4$ such branches is a tree of degrees $\le4$ whose root is the vertex of case (i). As isomorphisms
preserve the centroid, the classes of such trees are the classes of these rooted trees, which are the multisets of $4$ slots each
holding a ternary tree (possibly empty) with $<\frac n2$ vertices, of total size $n-1$. Proposition~\ref{prop:multisets}, with the empty tree
allowed as an object of size $0$ (the proof is unchanged, as $r_0=1$), gives the first term.

(ii) Cut the halving edge $uv$. The two halves, rooted at $u$ and $v$, are ternary rooted trees with $\frac n2$ vertices, and conversely
joining the roots of two such trees gives a tree of degrees $\le4$ with a halving edge. Isomorphisms may swap the halves, so the
classes correspond to unordered pairs, with repetition, of classes of ternary trees of size $\frac n2$: $\binom{r_{n/2}+1}2$ of them.
\end{proof}

For $n=6$: the branches at a centroid have at most $2$ vertices, and the multisets of $4$ such branches of total size $5$ are $\{2,2,1,0\}$
and $\{2,1,1,1\}$, that is, 3-methylpentane and 2,2-dimethylbutane; and $r_3=2$ gives $\binom32=3$ trees with a halving edge: hexane,
2-methylpentane and 2,3-dimethylbutane. Five hexanes, as in Figure~\ref{fig:hexanes}. The sequence goes on, $1$, $1$, $1$, $2$, $3$, $5$, $9$, $18$, $35$, $75$, $159$, $355$, $802$, $1858$, $4347$, \dots,
and the notebook checks it against a brute-force enumeration of trees. There are $366\,319$ isomers of $\mathrm C_{20}\mathrm H_{42}$ and more
than $4\times10^{9}$ of $\mathrm C_{30}\mathrm H_{62}$; Cayley's thirteen carbons took him a great deal of patience, and take a computer a millisecond.

\section{Gaussian binomial coefficients}

\begin{think}
How many lines through the origin are there in the plane $\mathbb F_2^2$? In $\mathbb F_3^2$? How many planes through the origin in $\mathbb F_2^3$?
\end{think}

In 1808 Gauss, studying the sign of the sums that now bear his name, needed polynomials in a variable $q$ that behave like binomial
coefficients. A century and a half later they turned out to count subspaces of vector spaces over finite fields, a fact used
every day in coding theory. For an integer $n\ge0$ let
\[
\begin{gathered}[n]_q=1+q+\dots+q^{n-1}=\frac{q^n-1}{q-1},\qq{}[n]_q!=[1]_q[2]_q\cdots[n]_q,\qq \binom nk_{\!q}=\frac{[n]_q!}{[k]_q!\,[n-k]_q!}\quad(0\le k\le n),\end{gathered}
\]
and $\binom nk_q=0$ if $k<0$ or $k>n$. At $q=1$ these are $n$, $n!$ and $\binom nk$.

\begin{proposition}[$q$-Pascal]\label{prop:qpascal}
For $n\ge1$,
\[
\begin{aligned}{}&\binom nk_{\!q}=\binom{n-1}{k-1}_{\!q}+q^k\binom{n-1}k_{\!q}\eqbr=q^{n-k}\binom{n-1}{k-1}_{\!q}+\binom{n-1}k_{\!q}.\end{aligned}
\]
Hence $\binom nk_q$ is a polynomial in $q$ with non-negative integer coefficients, of degree $k(n-k)$.
\end{proposition}
\begin{proof}
For $0<k<n$, divide the first identity by $\frac{[n-1]_q!}{[k]_q!\,[n-k]_q!}$: it becomes $[n]_q=[k]_q+q^k[n-k]_q$, which says $1+\dots+q^{n-1}=(1+\dots+q^{k-1})+(q^k+\dots+q^{n-1})$.
The second is the same with $k$ and $n-k$ exchanged, using the symmetry $\binom nk_q=\binom n{n-k}_q$. For $k=0$ or $k=n$ both sides are $1$. The
last claims follow by induction on $n$.
\end{proof}

\begin{theorem}\label{thm:subspaces}
Let $q$ be a prime power. The number of subspaces of dimension $k$ of $\mathbb F_q^n$ is $\binom nk_q$. The number of invertible $n\times n$
matrices over $\mathbb F_q$ is $|\mathrm{GL}_n(\mathbb F_q)|=(q^n-1)(q^n-q)\cdots(q^n-q^{n-1})$.
\end{theorem}
\begin{proof}
The linear algebra of Tome~0, Chapter~7 works word for word over any field, and $\mathbb F_q$ exists by Tome~2, Theorem~9.6. Count the
sequences $(v_1,\dots,v_k)$ of linearly independent vectors of $\mathbb F_q^n$: $v_1$ is any of the $q^n-1$ non-zero vectors, and $v_{i+1}$ is any vector
outside the span of $v_1,\dots,v_i$, which has $q^i$ elements: $(q^n-1)(q^n-q)\cdots(q^n-q^{k-1})$ sequences. For $k=n$ these are the columns of the
invertible matrices. Each sequence is a basis of exactly one subspace of dimension $k$, and each such subspace, being isomorphic
to $\mathbb F_q^k$, has $(q^k-1)(q^k-q)\cdots(q^k-q^{k-1})$ bases. So the number of subspaces is
\[
\begin{aligned}{}&\prod_{i=0}^{k-1}\frac{q^n-q^i}{q^k-q^i}=\prod_{i=0}^{k-1}\frac{q^{n-i}-1}{q^{k-i}-1}\eqbr=\frac{[n]_q[n-1]_q\cdots[n-k+1]_q}{[k]_q[k-1]_q\cdots[1]_q}=\binom nk_{\!q}.\end{aligned}
\]
\end{proof}

So $\mathbb F_2^3$ has $\binom31_2=1+2+4=7$ lines and $7$ planes: the Fano plane, which Chapter~12 will meet again. And the polynomial has
a meaning for every $q$, not only prime powers. An \emph{inversion} of a word $w\in\{0,1\}^n$ is a pair of positions $i<j$ with $w_i=1$ and $w_j=0$.

\begin{theorem}\label{thm:qinv}
$\displaystyle\binom nk_{\!q}=\sum_wq^{\operatorname{inv}(w)}$, the sum over the words $w\in\{0,1\}^n$ with $k$ ones.
\end{theorem}
\begin{proof}
Call the right-hand side $I(n,k)$; $I(n,0)=I(n,n)=1$. For $0<k<n$, split by the last letter. If $w=w'1$, the final $1$ is in no inversion,
so these words contribute $I(n-1,k-1)$. If $w=w'0$, the final $0$ makes an inversion with each of the $k$ ones of $w'$, so they contribute
$q^kI(n-1,k)$. So $I$ satisfies the first recursion of Proposition~\ref{prop:qpascal}, with the same initial values.
\end{proof}

Draw $w$ as a lattice path, a $1$ being a step right and a $0$ a step up: the inversions are the unit squares above the path
and inside the $k\times(n-k)$ rectangle. So $\binom nk_q$ is the generating function of lattice paths by the area they cut off, and its
coefficients are symmetric and unimodal (the unimodality is a deep theorem of Sylvester, 1878). Exercise~\ref{ex:rref} explains
why subspaces and inversions give the same count: row reduction turns each subspace into a lattice path.

\begin{theorem}[$q$-binomial theorem]\label{thm:qbinomial}
$\displaystyle\prod_{i=0}^{n-1}(1+q^ix)=\sum_{k=0}^nq^{k(k-1)/2}\binom nk_{\!q}\,x^k$.
\end{theorem}
\begin{proof}
By induction on $n$, the case $n=0$ being $1=1$. Multiplying the identity for $n-1$ by $1+q^{n-1}x$, the coefficient of $x^k$ on the left
becomes
\[
\begin{aligned}{}&q^{k(k-1)/2}\binom{n-1}k_{\!q}+q^{n-1}q^{(k-1)(k-2)/2}\binom{n-1}{k-1}_{\!q}\eqbr=q^{k(k-1)/2}\Big(\binom{n-1}k_{\!q}+q^{n-k}\binom{n-1}{k-1}_{\!q}\Big)\eqbr=q^{k(k-1)/2}\binom nk_{\!q},\end{aligned}
\]
since $n-1+\frac{(k-1)(k-2)}2-\frac{k(k-1)}2=n-k$, and by the second recursion of Proposition~\ref{prop:qpascal}.
\end{proof}

\begin{beauty}
At $q=1$ everything collapses to the binomial theorem and to subsets of a set. Jacques Tits joked in 1957 that a set is a vector
space over ``the field with one element'', and that combinatorics is linear algebra at $q=1$. The joke became a research
programme, still alive: a geometry over $\mathbb F_1$ in which, some hope, the Riemann hypothesis might become as provable as it is
for curves over finite fields.
\end{beauty}

\section{Exercises}

\begin{exercise}[$\star$]\label{ex:binomdouble}
Prove by double counting, without computation: (a) $\sum_kk\binom nk=n2^{n-1}$; (b) $\sum_k\binom nk^2=\binom{2n}n$; (c) $\sum_{k=r}^n\binom kr=\binom{n+1}{r+1}$.
\end{exercise}

\begin{exercise}[$\star\star$]\label{ex:partitionlattice}
The partitions of $\{1,\dots,n\}$ into blocks, ordered by refinement (a partition is below another if each of its blocks lies in a
block of the other), form the \emph{partition lattice} $\Pi_n$, with bottom $\hat0$ (all singletons) and top $\hat1$ (one block). Show that the
interval $[\hat0,\pi]$ is isomorphic to the product of the $\Pi_{|B|}$ over the blocks $B$ of $\pi$, compute $\mu(\hat0,\hat1)$ for $\Pi_3$ and $\Pi_4$, and guess
the general formula.
\end{exercise}

\begin{exercise}[$\star\star$]\label{ex:lyndon}
Order words over $\{0,\dots,q-1\}$ lexicographically. A \emph{Lyndon word} is a non-empty word strictly smaller than all its proper
rotations. Show that the number of Lyndon words of length $n$ is $\frac1n\sum_{d\mid n}\mu(n/d)\,q^d$, and list the binary Lyndon words of
lengths $4$ and $6$.
\end{exercise}

\begin{exercise}[$\star$]\label{ex:leaf}
Show that in a uniformly random tree on $\{1,\dots,n\}$, the probability that vertex $1$ is a leaf is $(1-\frac1n)^{n-2}$, and that the expected
number of leaves is $n(1-\frac1n)^{n-2}\sim\frac ne$.
\end{exercise}

\begin{exercise}[$\star\star$]\label{ex:pitman}
\emph{Pitman's proof.} Count in two ways the sequences of directed edges which, added one by one to the graph with vertices
$\{1,\dots,n\}$ and no edges, produce a rooted tree with all edges directed away from the root. Deduce Cayley's formula.
\end{exercise}

\begin{exercise}[$\star\star$]\label{ex:fusscatalan}
A plane tree is \emph{$r$-ary} if every vertex has $0$ or $r$ children. Show that the number of $r$-ary plane trees with $n$ vertices having
$r$ children is $\frac1{rn+1}\binom{rn+1}n=\frac1{(r-1)n+1}\binom{rn}n$ (the \emph{Fuss--Catalan numbers}). How many ternary trees have $3$ internal
vertices?
\end{exercise}

\begin{exercise}[$\star$]\label{ex:cube}
Show that the number of ways to colour the faces of a cube with $q$ colours, up to rotation, is $\frac1{24}(q^6+3q^4+12q^3+8q^2)$. How many
with $2$ colours?
\end{exercise}

\begin{exercise}[$\star\star$]\label{ex:graphs5}
Compute the cycle index of $S_5$ acting on the $10$ pairs of $\{1,\dots,5\}$, and deduce that there are $34$ graphs on $5$ unlabelled
vertices.
\end{exercise}

\begin{exercise}[$\star\star$]\label{ex:rootedrec}
From $T(x)=x\prod_n(1-x^n)^{-t_n}$, derive the recursion $nt_{n+1}=\sum_{m=1}^ns_m\,t_{n+1-m}$, where $s_m=\sum_{d\mid m}d\,t_d$, and compute $t_1,\dots,t_7$.
Draw the four rooted trees with $4$ vertices.
\end{exercise}

\begin{exercise}[$\star\star$]\label{ex:rref}
Every subspace $V$ of dimension $k$ of $\mathbb F_q^n$ is the row space of exactly one $k\times n$ matrix in \emph{reduced row echelon form}: each row
starts with a $1$ (its pivot), the pivots move strictly to the right from row to row, and the pivot columns are zero outside the
pivots. Prove this, and deduce a second proof of Theorem~\ref{thm:qinv} when $q$ is a prime power, then for every $q$.
\end{exercise}

\begin{exercise}[$\star$]\label{ex:glprob}
Show that a random $n\times n$ matrix over $\mathbb F_2$, with independent uniform entries, is invertible with probability $\prod_{i=1}^n(1-2^{-i})$, and
that this decreases to a limit between $\frac14$ and $\frac38$. (It is $0.2887\ldots$)
\end{exercise}

\begin{exercise}[$\star\star\star$]\label{ex:menage}
\emph{The m\'enage problem} (Lucas, 1891). In how many ways can $n$ couples sit at a round table, men and women alternating, with
nobody next to their partner? Seat the women first; show that the men can then be seated in
\[
M_n=\sum_{k=0}^n(-1)^k\,\frac{2n}{2n-k}\binom{2n-k}k\,(n-k)!
\]
ways (Touchard, 1934), and check $M_3=1$ and $M_4=2$. You may first show that a cycle of $m$ points has $\frac m{m-k}\binom{m-k}k$ sets of $k$
points with no two adjacent ($0\le k<m$).
\end{exercise}

\section{Hints}

\begin{hint}{ex:binomdouble}
(a) Committees with a chair. (b) Choose $n$ people from $n$ women and $n$ men. (c) Sort the $(r+1)$-subsets of $\{1,\dots,n+1\}$ by their largest
element.
\end{hint}
\begin{hint}{ex:partitionlattice}
A partition below $\pi$ is a choice of a partition of each block of $\pi$. Then use Proposition~\ref{prop:mobiusproduct}.
\end{hint}
\begin{hint}{ex:lyndon}
In each class of $n$ rotations of a primitive word, take the smallest. Is a Lyndon word primitive?
\end{hint}
\begin{hint}{ex:leaf}
Lemma~\ref{lem:pruferdeg}: vertex $1$ is a leaf if and only if its letter is absent from the code.
\end{hint}
\begin{hint}{ex:pitman}
When the current forest has $k$ rooted trees, which edges can be added?
\end{hint}
\begin{hint}{ex:fusscatalan}
Proposition~\ref{prop:lukasiewicz}, with numbers of children in $\{0,r\}$.
\end{hint}
\begin{hint}{ex:cube}
The rotations: the identity, $\pm90^\circ$ and $180^\circ$ about the $3$ axes through opposite faces, $\pm120^\circ$ about the $4$ diagonals, $180^\circ$ about the
$6$ axes through midpoints of opposite edges. Count the cycles of each on the faces.
\end{hint}
\begin{hint}{ex:graphs5}
There are seven cycle types in $S_5$. For each, follow the pairs as in the text.
\end{hint}
\begin{hint}{ex:rootedrec}
Take the logarithmic derivative: $x\frac{T'(x)}{T(x)}=1+\sum_nt_n\frac{nx^n}{1-x^n}$.
\end{hint}
\begin{hint}{ex:rref}
The pivot columns of $V$ are the positions $j$ at which some vector of $V$ has its first non-zero coordinate. The free entries of an
echelon matrix are inversions of the word marking the pivot columns.
\end{hint}
\begin{hint}{ex:glprob}
Theorem~\ref{thm:subspaces}, and $\prod(1-a_i)\ge1-\sum a_i$ for $a_i\in[0,1]$.
\end{hint}
\begin{hint}{ex:menage}
Man $i$ may not sit in two of the $n$ gaps. Use inclusion--exclusion over sets of forbidden (man, gap) pairs, and note which of them
can occur together.
\end{hint}

\section{Solutions}

\begin{solution}{ex:binomdouble}
\idea Each sum counts a set of structures sorted by a parameter; find the structures, and count them once without sorting.
\writeup (a) Count committees with a chair chosen from $n$ people. Choosing the committee first, of size $k$, then the chair: $\sum_kk\binom nk$.
Choosing the chair first, then any subset of the other $n-1$: $n2^{n-1}$. (b) Choose $n$ people from $n$ women and $n$ men: $\binom{2n}n$ ways. Sorting
by the number $k$ of women chosen, there are $\binom nk\binom n{n-k}=\binom nk^2$ choices with $k$ women. (c) An $(r+1)$-subset of $\{1,\dots,n+1\}$ with largest element
$k+1$ consists of $k+1$ and an $r$-subset of $\{1,\dots,k\}$: $\binom kr$ of them, for $r\le k\le n$.
\end{solution}

\begin{solution}{ex:partitionlattice}
\idea A refinement of $\pi$ refines each block independently, so the interval below $\pi$ is a product, and $\mu$ is multiplicative.
\writeup If $\pi$ has blocks $B_1,\dots,B_r$, a partition $\sigma\le\pi$ is the same thing as a partition $\sigma_i$ of each $B_i$, and $\sigma\le\sigma'\le\pi$ if and only if
$\sigma_i\le\sigma'_i$ for all $i$: so $[\hat0,\pi]\cong\Pi_{|B_1|}\times\dots\times\Pi_{|B_r|}$, and by Proposition~\ref{prop:mobiusproduct} $\mu(\hat0,\pi)=\prod_i\mu_{|B_i|}$, where
$\mu_m=\mu_{\Pi_m}(\hat0,\hat1)$ and $\mu_1=1$. In $\Pi_2$, $\mu_2=-1$. In $\Pi_3$, the three partitions of type $2+1$ have $\mu=-1$, so $\mu_3=-(1-3)=2$.
In $\Pi_4$, below $\hat1$ there are: $\hat0$ ($\mu=1$), six partitions of type $2+1+1$ ($\mu=-1$), three of type $2+2$ ($\mu=\mu_2^2=1$) and four of type
$3+1$ ($\mu=\mu_3=2$). So $\mu_4=-(1-6+3+8)=-6$. The pattern $1,-1,2,-6$ suggests $\mu_n=(-1)^{n-1}(n-1)!$, which is true (Stanley,
\emph{Enumerative Combinatorics}, Volume~1, Example~3.10.4); the notebook checks it up to $n=6$.
\end{solution}

\begin{solution}{ex:lyndon}
\idea Group primitive words into their rotation classes and pick the least element of each.
\writeup A rotation of a power $u^m$ is a power of a rotation of $u$, so rotations of primitive words are primitive, and a primitive
word $w$ of length $n$ has $n$ distinct rotations (the rotations fixing $w$ form the trivial subgroup, Lemma~\ref{lem:primitive}). So the
$M(n)$ primitive words split into $M(n)/n$ classes of $n$ distinct words, and each class has a unique smallest element, which is
strictly smaller than its other rotations: a Lyndon word. Conversely a Lyndon word is primitive, since for $w=u^m$ with $m\ge2$ the
rotation by $|u|$ equals $w$; and it is the smallest of its class. So Lyndon words of length $n$ correspond to classes, and their number
is $\frac{M(n)}n=\frac1n\sum_{d\mid n}\mu(n/d)q^d$. For $q=2$: length $4$ gives $\frac{16-4}4=3$, namely $0001$, $0011$, $0111$; length $6$ gives $\frac{64-8-4+2}6=9$, namely
$000001$, $000011$, $000101$, $000111$, $001011$, $001101$, $001111$, $010111$, $011111$.
\end{solution}

\begin{solution}{ex:leaf}
\idea Through the Pr\"ufer bijection, a uniform tree is a uniform word, and ``$1$ is a leaf'' becomes ``the letter $1$ is missing''.
\writeup By Lemma~\ref{lem:pruferdeg}, $1$ is a leaf if and only if $1$ does not occur in the code. The codes avoiding $1$ are the $(n-1)^{n-2}$ words
on the other letters, so the probability is $\frac{(n-1)^{n-2}}{n^{n-2}}=(1-\frac1n)^{n-2}$. By symmetry (or linearity of expectation, summing
indicators over the $n$ vertices), the expected number of leaves is $n(1-\frac1n)^{n-2}$, and $(1-\frac1n)^{n-2}=\exp((n-2)\ln(1-\frac1n))\to e^{-1}$.
\end{solution}

\begin{solution}{ex:pitman}
\idea Build the tree edge by edge: the result is a rooted tree with an ordering of its edges, or, counted forwards, a sequence
of choices whose number does not depend on the earlier ones.
\writeup Call such a sequence of $n-1$ edges a \emph{building sequence}. First count: choose a tree ($T_n$ ways, $T_n$ the number of trees on
$\{1,\dots,n\}$), a root ($n$ ways, which directs the edges away from it), and an order of its $n-1$ edges: $T_n\,n\,(n-1)!$ building sequences.
Second count: after $n-k$ edges have been added, the graph is a forest of $k$ rooted trees with edges directed away from the roots
(every vertex has at most one incoming edge, and there are no cycles). The edge $u\to v$ can be added, keeping this property, if and
only if $v$ is the root of a tree not containing $u$: $n$ choices of $u$ and then $k-1$ choices of $v$. So there are $\prod_{k=2}^nn(k-1)=n^{n-1}(n-1)!$
building sequences. Equating, $T_n=n^{n-2}$.
\end{solution}

\begin{solution}{ex:fusscatalan}
\idea An $r$-ary tree is determined by its \L ukasiewicz word, a word in the letters $0$ and $r$; count the arrangements and divide by
the length.
\writeup An $r$-ary tree with $n$ vertices having $r$ children has $rn$ children in all, hence $rn+1$ vertices, of which $rn+1-n$ have $0$ children.
By Proposition~\ref{prop:lukasiewicz}, the number of such plane trees is $\frac1{rn+1}$ times the number of arrangements of $n$ letters $r$ and
$(r-1)n+1$ letters $0$, that is $\frac1{rn+1}\binom{rn+1}n$. And $\frac1{rn+1}\binom{rn+1}n=\frac{(rn)!}{n!\,((r-1)n+1)!}=\frac1{(r-1)n+1}\binom{rn}n$. For $r=2$ these are the Catalan
numbers (binary trees); for $r=3$, $n=3$: $\frac1{10}\binom{10}3=12$ ternary trees.
\end{solution}

\begin{solution}{ex:cube}
\idea P\'olya with the faces as slots: list the rotations by type and count their cycles on the six faces.
\writeup The identity has $6$ cycles. A rotation by $\pm90^\circ$ about a face axis fixes the two faces on the axis and cycles the four others:
$3$ cycles, $6$ such rotations. A rotation by $180^\circ$ about a face axis fixes two faces and swaps the others in pairs: $4$ cycles, $3$
rotations. A rotation by $\pm120^\circ$ about a diagonal cycles the three faces at each end: $2$ cycles, $8$ rotations. A rotation by $180^\circ$
about an edge axis swaps the faces in three pairs: $3$ cycles, $6$ rotations. That is $1+6+3+8+6=24$ rotations, and Burnside gives
$\frac1{24}(q^6+6q^3+3q^4+8q^2+6q^3)=\frac1{24}(q^6+3q^4+12q^3+8q^2)$. For $q=2$: $\frac{64+48+96+32}{24}=10$.
\end{solution}

\begin{solution}{ex:graphs5}
\idea Follow the pairs under one permutation of each cycle type; a pair inside a cycle of $\sigma$ or across two cycles moves along a
predictable orbit.
\writeup The classes of $S_5$, with their sizes and the cycles induced on the $10$ pairs:
identity ($1$): $t_1^{10}$; $(12)$ ($10$): $\{1,2\}$ and the $3$ pairs inside $\{3,4,5\}$ are fixed, $\{1,j\}\leftrightarrow\{2,j\}$ for $j=3,4,5$: $t_1^4t_2^3$;
$(12)(34)$ ($15$): $\{1,2\},\{3,4\}$ fixed, $\{1,3\}\leftrightarrow\{2,4\}$, $\{1,4\}\leftrightarrow\{2,3\}$, $\{1,5\}\leftrightarrow\{2,5\}$, $\{3,5\}\leftrightarrow\{4,5\}$: $t_1^2t_2^4$;
$(123)$ ($20$): the pairs inside $\{1,2,3\}$, the pairs $\{i,4\}$ and the pairs $\{i,5\}$ form three $3$-cycles, $\{4,5\}$ is fixed: $t_1t_3^3$;
$(123)(45)$ ($20$): a $3$-cycle inside $\{1,2,3\}$, $\{4,5\}$ fixed, and the six pairs $\{i,j\}$ with $i\le3<j$ form one $6$-cycle: $t_1t_3t_6$;
$(1234)$ ($30$): the four sides of the square $1234$ form a $4$-cycle, the two diagonals a $2$-cycle, the pairs $\{i,5\}$ a $4$-cycle: $t_2t_4^2$;
$(12345)$ ($24$): two $5$-cycles (sides and diagonals of the pentagon): $t_5^2$. The sizes add up to $120$. With $t_j=2$:
\[
\begin{aligned}{}&\tfrac1{120}\big(2^{10}+10\cdot2^7+15\cdot2^6+20\cdot2^4\eqbr{}+20\cdot2^3+30\cdot2^3+24\cdot2^2\big)\eqbr=\tfrac{1024+1280+960+320+160+240+96}{120}\eqbr=\tfrac{4080}{120}=34.\end{aligned}
\]
\end{solution}

\begin{solution}{ex:rootedrec}
\idea A product becomes a sum under the logarithmic derivative, and each factor $(1-x^n)^{-t_n}$ contributes a geometric series.
\writeup The logarithmic derivative $F\mapsto xF'/F$ turns products into sums, by the product rule (for formal power series with
non-zero constant term, Chapter~4; each coefficient involves finitely many factors). For $F=x$ it is $1$; for $F=(1-x^n)^{-t_n}$ it is
$t_n\frac{nx^n}{1-x^n}=t_n\sum_{j\ge1}nx^{nj}$. Collecting the terms in $x^m$: $\frac{xT'}T=1+\sum_{m\ge1}s_mx^m$. Multiply by $T$ and compare the coefficients of $x^{n+1}$:
$(n+1)t_{n+1}=t_{n+1}+\sum_{m=1}^ns_mt_{n+1-m}$, the recursion. With $t_1=1$: $s_1=1$, $t_2=1$; $s_2=3$, $2t_3=1+3$, $t_3=2$; $s_3=7$, $3t_4=2+3+7$, $t_4=4$;
$s_4=19$, $4t_5=4+6+7+19$, $t_5=9$; $s_5=46$, $5t_6=9+12+14+19+46$, $t_6=20$; $s_6=129$, $6t_7=20+27+28+38+46+129$, $t_7=48$. The four rooted trees
with $4$ vertices: a path rooted at an end; a root with one child having two children; a root with two children, one of which has
a child; a root with three children.
\end{solution}

\begin{solution}{ex:rref}
\idea The pivot columns of a subspace can be read without choosing a basis; once they are known, the echelon rows are forced.
\writeup Let $J$ be the set of positions $j$ such that some vector of $V$ has its first non-zero coordinate at $j$. For each $j\in J$ choose such a
vector; these vectors are independent (their leading positions differ), and a vector of $V$ not in their span could be reduced,
by subtracting multiples of them, to a non-zero vector of $V$ whose leading position is not in $J$, absurd. So $|J|=k$. Reducing
each chosen vector by the others (Gaussian elimination) gives a basis $u_j$, $j\in J$, with $u_j$ leading at $j$ with coefficient $1$ and zero
at the other positions of $J$: the rows of a reduced echelon matrix. It is unique: if $u'_j$ is another such basis, $u_j-u'_j\in V$ vanishes
at every position of $J$, so its leading position is not in $J$, and $u_j=u'_j$. Conversely the row space of a reduced echelon matrix has
dimension $k$. So the subspaces correspond to the reduced echelon matrices. Encode the pivot columns $J$ by the word $w$ with $w_j=1$ for
$j\in J$. The entries of row $i$ that are not forced are those in the non-pivot columns to the right of its pivot: summing over rows,
there are as many free entries as pairs (a $1$ of $w$, a later $0$ of $w$), that is, $\operatorname{inv}(w)$. So the number of subspaces is $\sum_wq^{\operatorname{inv}(w)}$, which with
Theorem~\ref{thm:subspaces} proves Theorem~\ref{thm:qinv} for all prime powers $q$. Two polynomials that agree at infinitely many
points are equal (Tome~0, Chapter~3), so the identity holds for every $q$.
\end{solution}

\begin{solution}{ex:glprob}
\idea Divide the number of invertible matrices by the number of all matrices; the product converges because its factors approach $1$
geometrically fast.
\writeup By Theorem~\ref{thm:subspaces}, the probability is $\frac{\prod_{i=0}^{n-1}(2^n-2^i)}{2^{n^2}}=\prod_{i=0}^{n-1}(1-2^{i-n})=\prod_{i=1}^n(1-2^{-i})$. Each new factor is $<1$, so the
products decrease to a limit $L$. Upper bound: $L\le(1-\frac12)(1-\frac14)=\frac38$. Lower bound: $\prod(1-a_i)\ge1-\sum a_i$ for $a_i\in[0,1]$, by induction, since
$(1-s)(1-a)\ge1-s-a$. So $L\ge\frac12(1-\sum_{i\ge2}2^{-i})=\frac12\cdot\frac12=\frac14$. Random linear codes and random linear hashing (Tome~C4) rely on such numbers:
a random square binary matrix is singular more often than not.
\end{solution}

\begin{solution}{ex:menage}
\idea Inclusion--exclusion over forbidden seats; the forbidden (man, gap) pairs form a cycle in which only non-adjacent pairs can
occur together, so the count reduces to independent sets on a cycle.
\writeup \emph{Cycles.} On a path of $p$ points, sets of $k$ points with no two adjacent number $\binom{p-k+1}k$: remove one point after each chosen
point but the last, getting any $k$-subset of $p-k+1$ points. On a cycle of $m$ points $1,\dots,m$: the sets avoiding point $1$ live on a
path of $m-1$ points, $\binom{m-k}k$ of them; the sets containing $1$ avoid $2$ and $m$, and the rest lives on a path of $m-3$ points, $\binom{m-k-1}{k-1}$ of them.
The total is $\binom{m-k}k+\frac k{m-k}\binom{m-k}k=\frac m{m-k}\binom{m-k}k$.

\emph{M\'enage.} With the women seated, the men go in the $n$ gaps, gap $j$ lying between women $j$ and $j+1$ (mod $n$), and man $i$ must avoid
gaps $i-1$ and $i$. A seating is a permutation $\sigma$ (man $i$ to gap $\sigma(i)$) avoiding the $2n$ forbidden cells $(i,i-1),(i,i)$. List them cyclically:
$(1,n),(1,1),(2,1),(2,2),\dots,(n,n-1),(n,n)$; two of them share a man or a gap exactly when they are consecutive in this cyclic list (for
$n\ge3$; check $n\le2$ by hand). By inclusion--exclusion, the number of permutations using no forbidden cell is $\sum_S(-1)^{|S|}N(S)$ over
sets $S$ of forbidden cells, where $N(S)$ is the number of permutations using all cells of $S$. $N(S)=0$ if two cells of $S$ share a man or a
gap, and $N(S)=(n-|S|)!$ otherwise. The sets of the second kind with $|S|=k$ are the sets of $k$ pairwise non-adjacent points on a cycle of
$2n$ points: $\frac{2n}{2n-k}\binom{2n-k}k$ (for $k=n$ it is $2$, the two alternate halves). This gives Touchard's formula. For $n=3$:
$6-6\cdot2+9\cdot1-2\cdot1=1$. For $n=4$: $24-8\cdot6+20\cdot2-16\cdot1+2\cdot1=2$. The answer to Lucas's question is $2\cdot n!\cdot M_n$ (seat the women in one of
$2\,n!$ ways: choose the alternate seats and the order of the women).
\end{solution}

\section*{Sources and further reading}

\begin{itemize}
\item A.~Cayley, \emph{On the theory of the analytical forms called trees}, Philosophical Magazine 13 (1857), 172--176;
\emph{On the analytical forms called trees, with application to the theory of chemical combinations}, Report of the British
Association for the Advancement of Science (1875), 257--305; \emph{A theorem on trees}, Quarterly Journal of Pure and Applied
Mathematics 23 (1889), 376--378.
\item C.~W.~Borchardt, \emph{\"Uber eine der Interpolation entsprechende Darstellung der Eliminations-Resultante}, Journal f\"ur
die reine und angewandte Mathematik 57 (1860), 111--121.
\item H.~Pr\"ufer, \emph{Neuer Beweis eines Satzes \"uber Permutationen}, Archiv der Mathematik und Physik 27 (1918), 142--144.
\item J.~Pitman, \emph{Coalescent random forests}, Journal of Combinatorial Theory A 85 (1999), 165--193. A research paper
whose first pages contain the proof of Exercise~\ref{ex:pitman}, then use the same building process to study random forests
and the coalescence of clusters, a model of physics and population genetics.
\item G.-C.~Rota, \emph{On the foundations of combinatorial theory I. Theory of M\"obius functions}, Zeitschrift f\"ur
Wahrscheinlichkeitstheorie 2 (1964), 340--368; P.~Hall, \emph{The Eulerian functions of a group}, Quarterly Journal of
Mathematics 7 (1936), 134--151.
\item A.~Dvoretzky and Th.~Motzkin, \emph{A problem of arrangements}, Duke Mathematical Journal 14 (1947), 305--313;
G.~N.~Raney, \emph{Functional composition patterns and power series reversion}, Transactions of the AMS 94 (1960), 441--451.
\item C.~Jordan, \emph{Sur les assemblages de lignes}, Journal f\"ur die reine und angewandte Mathematik 70 (1869), 185--190.
\item G.~P\'olya, \emph{Kombinatorische Anzahlbestimmungen f\"ur Gruppen, Graphen und chemische Verbindungen}, Acta
Mathematica 68 (1937), 145--254; English translation with commentary by R.~C.~Read, \emph{Combinatorial Enumeration of Groups,
Graphs, and Chemical Compounds}, Springer, 1987.
\item H.~R.~Henze and C.~M.~Blair, \emph{The number of structurally isomeric alcohols of the methanol series}, Journal of the
American Chemical Society 53 (1931), 3077--3085.
\item J.~Touchard, \emph{Sur un probl\`eme de permutations}, Comptes Rendus de l'Acad\'emie des Sciences 198 (1934), 631--633.
\item R.~P.~Stanley, \emph{Enumerative Combinatorics}, Volume~1, 2nd ed., Cambridge University Press, 2012 (M\"obius functions,
$q$-analogues), and \emph{Catalan Numbers}, Cambridge University Press, 2015.
\item F.~Harary and E.~M.~Palmer, \emph{Graphical Enumeration}, Academic Press, 1973. Counting graphs and trees with P\'olya's
theorem.
\item M.~Lothaire, \emph{Combinatorics on Words}, Addison-Wesley, 1983. Primitive words, Lyndon words and their factorisations.
\item M.~Aigner and G.~M.~Ziegler, \emph{Proofs from THE BOOK}, 6th ed., Springer, 2018. Four proofs of Cayley's formula.
\end{itemize}
